\documentclass[11pt, a4paper]{article}

% --- Pakker for font, språk og matematikk ---
\usepackage[utf8]{inputenc}
\usepackage[norsk]{babel}
\usepackage{mathpazo} % Palatino font
\usepackage{amsmath}
\usepackage{amssymb}
\usepackage{geometry}
\usepackage{hyperref}

\geometry{margin=2.5cm}

\title{Sannsynlighetskinematikk i Lukkede Multimodale Systemer: \\ 
Taktbasert Gjenkjenning i OMR og Lyd}
\author{Arbeidsnotat / Konseptskisse}
\date{\today}

\begin{document}

\maketitle

\begin{abstract}
Dette notatet skisserer en arkitektur for Optisk Musikkgjenkjenning (OMR) og multimodal audio-alignment. Ved å modellere musikknotasjon som et strengt regelstyrt, lukket system, kan vi utnytte Richard Jeffreys sannsynlighetskinematikk for å håndtere usikkerhet fra syns- og lydmodeller. Takten (eng. \textit{measure}) fungerer som den sentrale hierarkiske broen, en perfekt partisjon som tillater lokal konvergens av usikkerhet før globale grafer koordinerer relasjoner på tvers av taktstreker.
\end{abstract}

\section{Introduksjon}
Klassiske Bayesianske nettverk krever ofte absolutt sikkerhet (modus ponens) for å oppdatere betingede sannsynligheter. I analyse av historiske partiturer genererer maskinsyn og lydanalyse gjerne "skitne" observasjoner med betydelig usikkerhet. I stedet for å la denne usikkerheten forplante seg ukontrollert, formaliserer vi notasjonens iboende logikk som \textit{stive betingelser}, der usikkerhet i observasjonen kun oppdaterer vektingen av de underliggende, gjensidig utelukkende mulighetene.

\section{Sannsynlighetskinematikk (Jeffrey-betinging)}
Gitt et lukket utfallsrom $\Omega$ definert av takten, lar vi $\{B_1, B_2, \dots, B_n\}$ utgjøre en komplett partisjon av lokalt mulige grafiske eller musikalske tilstander (f.eks. at et symbol er en notehals, en taktstrek, eller støy). Partisjonen er gjensidig utelukkende og uttømmende: $\sum p(B_i) = 1$.

En syns- eller lydmodell observerer systemet og genererer en ny sannsynlighetsfordeling $q(B_i)$ for disse grunntilstandene. For en gitt musikalsk hypotese $x$, postulerer sannsynlighetskinematikken den kinematiske rigiditetsbetingelsen:
\begin{equation}
    q(x|B_i) = p(x|B_i)
\end{equation}
Dette betyr at de musikalske reglene er uavhengige av observasjonsfeil. Den nye marginale sannsynligheten for $x$ beregnes da via den generaliserte loven om total sannsynlighet:
\begin{equation}
    q(x) = \sum_{i=1}^{n} p(x|B_i)q(B_i)
\end{equation}
Operasjonelt lar dette oss isolere den visuelle støyen i vektoren $\mathbf{q}$, mens musikkteorien bevares i en overføringsmatrise $\mathbf{P}_x$.

\section{Takten som Avsnitt: Den Perfekte Partisjon}
Vi modellerer takten som en lokal, matematisk streng avgrensning. Analogt med tekstanalysens dokument-/avsnitt-hierarki, utgjør takten et lokalt mulighetsrom der summen av rytmiske verdier må samsvare nøyaktig med taktartens definisjon (f.eks. $\frac{4}{4}$).

Dette skaper en global rigid betingelse for den lokale partisjonen. Hvis $q(\text{lengde})$ for en gitt note er usikker, vil lengden på de sikrere elementene i samme takt \textit{tvinge} frem en løsning. Det lukkede rommet eliminerer matematisk umulige løsninger.

\section{Multimodal Alignment: Audio som Epistemisk Fasit}
Når systemet utvides med auditiv informasjon (lydopptak knyttet til partituret), forsterkes taktens funksjon som bindeledd. Lyd gir ekstremt presise data for tid (onsets/offsets) og frekvens (tonehøyde), selv der notebildet er uleselig.

Vi lar lydmodellen produsere en uavhengig observasjonsvektor $\mathbf{q}_{\text{audio}}$. Siden lyd forløper lineært i tid, fungerer taktslaget som synkroniseringspunktet (\textit{alignment}). Om synsmodellen er usikker på en akkord i Takt 4, vil $\mathbf{q}_{\text{audio}}$ levere en dominant sannsynlighet for de spesifikke frekvensene på det nøyaktige tidsavsnittet. Via Jeffrey-betinging kollapser dette den grafiske usikkerheten i takten, fordi logikken krever at det visuelle symbolet må produsere den auditive frekvensen.

\section{Koordinasjon på tvers av takter}
Mens takten sikrer lokal kinematikk, håndteres fenomener som krysser taktgrenser -- primært bindebuer (\textit{ties}) og fraseringsbuer (\textit{slurs}) -- via en to-pass grafarkitektur.

\begin{enumerate}
    \item \textbf{Pass 1 (Lokal Konvergens):} Ligningene løses lukket innad i hver takt. Buer detekteres med en lokal sannsynlighet (f.eks. \textit{Start\_Bue} i Takt 1), men påvirker ikke taktens interne rytmiske sum.
    \item \textbf{Pass 2 (Global Graf):} Buer instansieres som rettede relasjoner (edges) mellom noder på tvers av takter. En bindebue fungerer som en rigid logisk operator som krever lik tonehøyde mellom nodene.
\end{enumerate}

Denne overføringen av sikkerhet på tvers av systemet tvinger frem en global konsensus, drevet av lokal matrisemultiplikasjon.

\end{document}